%% ch13 --- The AI engineer's kit: pydantic, httpx and the shape of an SDK
%%
%% Written September 2026. Every claim about an SDK in this chapter is
%% printed by code/ch13/sdk_shape.py out of the installed package, and every
%% number comes from code/measure/e08_validation.py. Nothing here was
%% written from memory; where the brief turned out to be wrong, CLAUDE.md's
%% *Resolved questions* says so.

\chapter{The AI engineer's kit: pydantic, httpx and the shape of an SDK}
\label{ch:aitoolkit}

\section{Three things in a trench coat}
\label{sec:ai-three-things}
\index{SDK!what one is made of}

You have added \pkg{openai} or \pkg{anthropic} to a project and the editor is
offering you a hundred names. Arriving from .NET, the reasonable expectation
is a vendor SDK: a generated client with its own object model, its own
options type, its own handler pipeline, and a layer of abstraction between
you and the wire that you will eventually have to fight.

That expectation is wrong, and it is wrong in the direction that makes your
life easier. Before you read on, answer this: a call to the SDK hands you
back an object. Not what is inside it \dash{} what \emph{class} is it an
instance of?

The answer is a pydantic model. The thing that fetched it is an httpx
client. The streaming form of the same call is an iterator. That is the
whole of it: three pieces of Python you can already read, assembled in the
way you would have assembled them.

\mermaidfig{sdk-anatomy}{What every model SDK in Python is made of. Each of
the three is a library you can read, and the SDK is the wiring between
them.}{the three parts of a model SDK}

\begin{csbox}
\code{HttpClient} plus \code{System.Text.Json} plus
\code{IAsyncEnumerable<T>}, and the SDK is the code between them. What
Python does not have is \code{IHttpClientFactory}: there is no container
handing you a pooled client, so the lifetime is yours to decide, and
\S\ref{sec:ai-which-httpx} is where that decision gets made.
\end{csbox}

This chapter stops at one model call and one structured output. Agents,
tools, middleware and graphs are a different book \dash{} they are the
subject of the companion volume, and re-teaching them here would give you
two accounts of one thing that will drift apart. What you get instead is
the floor those things stand on, which is the part that does not change
when the framework does.

\section{A structured output is a class you already wrote}
\label{sec:ai-structured}
\index{pydantic!structured output}

A model returns text. A \emph{structured output} is that text parsed into a
type you declared, and in Python the type is a pydantic model. You met
pydantic at a service boundary in Chapter~\ref{ch:services}; this is the
same mechanism pointed the other way, at a reply instead of a request.

\pyregion{code/ch13/structured.py}{models}{Two declarations of one
shape. The second differs by one line, and \S\ref{sec:ai-what-costs}
measures what that line costs.}{lst:ch13-models}

Two calls do all the work, and between them they are what an SDK's
structured-output helper is.

\pyregion{code/ch13/structured.py}{schema}{Out to the wire: the model
becomes a JSON schema, and the SDK adds what the provider wants on top of
it.}{lst:ch13-schema}

\mermaidfig{schema-round-trip}{One class, used twice: once to say what the
reply must look like, and once to check that it did.}{a structured output,
out and back}

Going out, \api{model\_json\_schema()} turns your class into the JSON
schema the request carries. Coming back,
\api{model\_validate\_json()}\index{model\_validate\_json@\texttt{model\_validate\_json}}
turns the reply's text into an instance of that class. The SDK's
contribution to the schema is two keys \dash{} a name, and
\code{additionalProperties: false} for a provider that promises to return
your fields and nothing else. Everything else on the wire is pydantic's.

That is worth sitting with, because it is the answer to a question that
otherwise looks deep. There is no registry, no code generation step and no
mapping layer. The shape is written down in one place, in the class, and
both halves of the round trip read it from there.

\section{What gets through}
\label{sec:ai-what-gets-through}
\index{pydantic!strict and lax}

Now the part where your .NET habits are actively unhelpful. You declared
\code{severity: int}. The provider, or a proxy in front of it, sends
\code{"2"} \dash{} the right value, as a string. Before you read on: what
does \api{model\_validate\_json()} do with it?

\pyregion{code/ch13/structured.py}{validate}{The same payload, twice. One
of these is what \code{JsonSerializer.Deserialize} would have
done.}{lst:ch13-validate}

\transcript{ch13-structured}

It coerces. \code{"2"} becomes \num{2}, \code{"true"} becomes \code{True},
and nothing anywhere is told. That is pydantic's default and it is called
\emph{lax} mode.

\begin{trapbox}
\code{System.Text.Json} refuses a JSON string for a numeric property unless
you opt in with \code{NumberHandling}. pydantic is the other way round:
coercion is the default and refusing is the option. So the reflex that
\emph{deserialisation validates} holds in C\# and not here \dash{} a
boundary written from that reflex accepts a contract change silently, and
the first thing you learn about it is a wrong number somewhere downstream.
\end{trapbox}

\code{strict=True} on the model's config buys back the behaviour you
expected. The useful question is not which is correct in the abstract but
what each one lets past, and that is countable.

\pyregion{code/ch13/what_survives.py}{plain}{The third strategy: the type a
C\# engineer writes first.}{lst:ch13-plain}

\pyfile{code/ch13/what_survives.py}{Seven replies that are valid JSON and
wrong. The corpus is imported by experiment E8 rather than copied, so this
table and that measurement cannot come apart.}{lst:ch13-survives}

\transcript{ch13-survives}

Read the bottom row first: of \val{e08.corpus} broken replies, lax refuses
\val{e08.reject.lax}, strict refuses \val{e08.reject.strict}, and the
dataclass refuses \val{e08.reject.plain}. The dataclass lets
\val{e08.silent.plain} of them through, every one of them wrong, because a
dataclass generates an \api{\_\_init\_\_} and checks not one of its
annotations at run time.

Then read the last line of the table, which is the one worth the price of
the chapter. The single payload the dataclass rejects is the one that was
\emph{fine}: a provider added a field. Both SDKs expect that to happen
\dash{} their own models are configured to keep unknown fields rather than
drop them, which \S\ref{sec:ai-which-httpx} prints \dash{} so the naive
strategy scores one refusal and it is a false alarm.

\section{What it costs, measured}
\label{sec:ai-what-costs}
\index{measurement!E8, validation cost}

The obvious next question is what that safety costs, and the obvious
expectation is a trade: validation is work, work takes time, pick a point
on the curve. Experiment E8 times the three strategies on the corpus above
and the expectation does not survive it.

Two results, on the machine this was written on. \textbf{Strict costs
nothing over lax} \dash{} the two came out within a few per cent of each
other, either side, run to run. And \textbf{pydantic's parse-and-validate
beats the standard library's parser alone}: on a reply of
\val{e08.fields.wide} fields, \api{json.loads} on its own took more than
\val{e08.json.factor} times as long as \api{model\_validate\_json} took to
parse the same bytes \emph{and} check every field against a type.

The mechanism is not mysterious once stated. \api{json.loads} builds a
Python dictionary of every key and value the provider sent, and then you
throw most of it away. pydantic's core parses the bytes in compiled code
straight into the model's fields, and never materialises the intermediate
dictionary at all. The wider the reply, the further ahead it gets \dash{}
which is the opposite of the shape you would draw if you thought of
validation as a tax levied on top of parsing.

\begin{note}
Neither of those two figures is on this page, and that is deliberate. A
microsecond is a property of a machine, so it cannot be committed to a book
that re-runs its own measurements and fails when one moves. What E8 commits
instead is a pair of ceilings, chosen generously and \emph{asserted on every
run}: strict may cost at most \val{e08.strict.overhead.pct} per cent more
than lax, and \api{json.loads} alone must be at least
\val{e08.json.factor} times slower. Each is taken as the best of
\val{e08.repeats} runs of \val{e08.iterations} calls, because the fastest
run is the one least disturbed by the scheduler where an average is an
average of the disturbances. If a later pydantic breaks either, the
build goes red and this section is wrong the same day. The counts in
\S\ref{sec:ai-what-gets-through} are not bounds: they are properties of the
code and identical on every machine.
\end{note}

So the trade the section opened with does not exist, and the advice is
unusually simple. Turn strict on at every boundary a model's reply crosses.
It is one line, it refuses two more classes of broken reply, and it is free.

\section{Which httpx}
\label{sec:ai-which-httpx}
\index{httpx!which distribution}

Everything above is about the pydantic leg. The transport leg has a trap in
it that no documentation will warn you about, and the only way to find it
is to read the package you installed.

So read it. The client inside the SDK is an httpx client, and this book pins
\pkg{httpx}. Before you read on: what does
\api{isinstance(client.\_client, httpx.Client)} return?

\pyfile{code/ch13/sdk_shape.py}{Every claim this chapter makes about an SDK,
printed out of the installed package. The assertions at the bottom mean CI
fails the day one of them stops being true.}{lst:ch13-sdk-shape}

\transcript{ch13-sdk-shape}

\pkg{httpx} and \pkg{httpx2} are two distributions. This book pins the
first, at version~\httpxver, because that is what you write a client with;
both SDKs depend on the second, and the client inside them is an
\pkg{httpx2} one. They share an ancestry and a spelling, and they are not
the same package, so \api{isinstance(client.\_client, httpx.Client)} is
\code{False} on an SDK client that is obviously an httpx client.

\begin{warning}
The two SDKs disagree about whether handing them the wrong one is an error.
Measured, not inferred: \pkg{anthropic} refuses an \api{httpx.Client} with a
\code{TypeError} that names both packages, and \pkg{openai} accepts it and
works \dash{} at least on a plain request, which is as far as this was
tested. \enquote{It worked when I tried it} is therefore not evidence that
the client you passed is the one the SDK expects.
\end{warning}

Two more things fall out of that listing. The SDKs' own models are
configured to \emph{keep} fields nobody declared rather than ignore them,
which is what makes a provider adding a field a non-event. And the default
timeouts are not a mistake you want to inherit: the read timeout the listing
prints is ten minutes, which is a sensible default for a provider that streams a long
answer and a terrible one for a request sitting inside an HTTP handler that
somebody is waiting on. Chapter~\ref{ch:services} gives that request a
budget; the SDK will not give it one for you.

The client's lifetime is the other decision. An httpx client holds a
connection pool, so one per process is right and one per request throws the
pool away every time \dash{} the same argument, and the same failure mode,
as the \api{AsyncClient} lifetime in Chapter~\ref{ch:asyncio}.

\section{Tokens are quantities, and money is not a float}
\label{sec:ai-tokens}
\index{cost!tokens as a quantity}

Every reply carries the counts that decide the bill, in a \api{usage} field
that is another pydantic model. A service that does not add them up learns
what it spent from the invoice, and by then the run that spent it is a
fortnight old.

\begin{csbox}
A price is money, so it is a \api{Decimal} and never a \code{float}, for
exactly the reason it is \code{decimal} and never \code{double} in C\#.
The difference is that the C\# compiler makes you say which one you meant
and Python does not: an untyped literal is a float, \code{0.1 + 0.2} is not
\code{0.3} here either, and a month of per-call costs summed as floats
drifts. \api{Decimal} also takes its precision from a context rather than
from the type, so quantise deliberately instead of trusting a default.
\end{csbox}

Exercise~\ref{ex:ch13-cost} is that arithmetic, done once, properly.

\section{Notebooks against scripts}
\label{sec:ai-notebooks}
\index{notebooks!and why this book has none}

Every tutorial you will meet for this material is a notebook, and nothing in
this book is one. That is a position, so here is the argument for it.

A notebook is not a file that runs. It is a REPL that remembers every cell
you have run, in the order you ran them, which is not the order they appear
on screen. Delete a cell and the name it defined is still bound. Run them
top to bottom on a fresh kernel and you get a different result from the one
on your screen, and the gap between those two is where the afternoon goes.
Nothing in this book could be a notebook, because every listing here is
executed on the pinned versions on every push, and \enquote{it worked in my
kernel} is not a thing CI can check.

A notebook is right for the work it was built for: looking at data, trying
something, keeping a figure beside the code that drew it. Use one. The rule
worth holding is narrower than \emph{avoid notebooks} \dash{} it is that
nothing a service imports lives in one. The moment a function is going to be
called by something other than you, it belongs in a module, with a test
beside it.

\section{One call, typed at both ends}
\label{sec:ai-call-model}

Here is everything above, assembled into the smallest useful thing:
one call with a type going in, a type coming out, a time budget, and a
record of what happened.

\pyfile{code/ch13/call_model.py}{The payoff. Nothing here is specific to a
provider, and it drops into the Chapter~\ref{ch:services} service as
it stands.}{lst:ch13-call-model}

\transcript{ch13-call-model}

Three details earn their place. The timeout is an argument with a default
measured in seconds rather than the SDK's ten minutes. The failure is
raised as one exception type of your own, so a caller handles
\emph{the model did not answer} rather than learning the provider's error
taxonomy \dash{} and it is named \code{...Error} and not
\code{...Exception}, which is the Python convention and which the linter
enforces. And the reply is checked for having parsed at all before it is
returned, because \api{output\_parsed} is optional and a caller that
ignores that gets a \code{None} two stack frames away from the cause.

\mermaidfig{two-boundaries}{The two edges of a service, and the same
mechanism at both.}{validation at both boundaries}

The listing runs, in CI, on every push, against no provider and no API key.
A local server that answers in the shape the SDK parses is indistinguishable
from a provider as far as the SDK is concerned, which is what makes a
chapter about model SDKs testable at all.

\begin{projectbox}
\textbf{Stage 02 \dash{} recording a model call.} The trace model from
Chapter~\ref{ch:testing} left one thing open: what a model call puts in an
event's \code{tags}. This chapter settles it. A call is \emph{two} events
and not one, because a call that never came back is the case most worth
being able to find, and a single event written after the reply cannot
represent it. The \code{tags} carry the token counts and whether the reply
parsed; they do not carry the prompt or the completion, because a trace is
kept, shipped and read by people who were not in the room.
\end{projectbox}

\section{Exercises}
\label{sec:ai-exercises}

\begin{exercise}{e13_01_strict_boundary}{A boundary that does not guess}
\code{Reply} coerces, so the string \code{"2"} arrives as an integer and
nobody is told. Make the parse refuse anything whose types are not exactly
what the class declares, and return \code{None} instead. Do not write the
checks by hand.
\end{exercise}

\begin{exercise}{e13_02_wire_schema}{The schema a provider will accept}
A provider that promises to return your fields and nothing else wants
\code{additionalProperties: false} on every object in the schema. A nested
model is a second object, and pydantic does not nest it where you would
expect \dash{} it lifts it out. Close all of them.
\end{exercise}

\begin{exercise}{e13_03_what_it_cost}{What one call cost}
\label{ex:ch13-cost}
Price a reply from its two token counts and the provider's two rates. The
answer is money, so it is not a float; and it is quantised, so that a month
of them adds up to the same number twice.
\end{exercise}

\begin{exercise}{e13_04_read_the_trace}{Why a call is two events}
Write the two questions a trace exists to answer: which models were called
and never answered, and what the whole run spent. Neither may assume the
trace is well formed \dash{} that is the point of having it.
\end{exercise}

\section{Where this stops}
\label{sec:ai-stops}

You can now read an SDK rather than invoke it. When the next one appears,
with a different resource name and a different envelope, the three questions
that settle it are: what pydantic model comes back, which HTTP client is
underneath, and what the streaming form of the call yields. The answers are
in the package you installed and nowhere else \dash{} not in a blog post,
not in a model's recollection, and not reliably in the documentation.

What is deliberately not here is everything built \emph{on} this: tool
calling, agent loops, retries with judgement in them, graphs, and the state
that goes with them. That is the companion volume's subject, and it starts
where this chapter ends. Chapter~\ref{ch:traceassert} takes the other road,
finishing the instrument this chapter taught to record a call, and shipping
it.
